
\chapter{Limitations and Threats to Validity}
\label{ch:limitations}

\section{Scope of the claims}

Every claim in this report is about one curated dataset, one field dataset, four
architectures and the specific protocol recorded in Table~\ref{tab:protocol}. Nothing
here establishes behaviour on other crops, other regions, other cameras or other
diseases.

\section{Dataset limitations}

PlantDoc test contains 236 images across 27 scored classes. A single accuracy
estimate on that sample has an interval roughly ten points wide. Paired designs and
multiple seeds reduce the damage, and they do not make the sample large.

Twelve mappings rest on a relationship weaker than ``same disease''. Three carry an
explicit caveat in the mapping file. If any of those is wrong, the field number moves
in a direction we cannot determine in advance.

Non-English and low-resolution field photographs are present in PlantDoc and are not
treated separately, so the field number averages over image quality bands.

\section{Statistical limitations}

The unit of independence for between-run comparisons is the seed, and there are
three. A paired interval with two degrees of freedom is very wide, which is why the
augmentation result is reported as no detectable effect rather than as a small
positive effect.

Bootstrap intervals over test images describe sampling variation in the test set.
They do not describe variation in training, and they are never used to make a claim
about model stability.

Calibration error at 15 bins on 236 images is itself noisy. The factor-of-223 rise
is far outside that noise, but individual bin values should not be read closely.

\section{Disclosed defects}

Two problems in the original pipeline are recorded here rather than removed from the
record.

\subsection{Protocol v1 could not reproduce the reference result}

The first training protocol failed to reproduce the reference accuracy. It is marked
as a blocker in the issue log. Every number in this report comes from the frozen
protocol that replaced it, and no result from protocol v1 is reported.

\subsection{SmallCNN does not replicate its earlier figure}

An early run recorded 96.01\,\% for SmallCNN. Under the frozen protocol and full
evaluation it reaches 91.60\,\%. The earlier figure is not used anywhere in this
report. The discrepancy is documented because a reader comparing this report to an
earlier draft would otherwise find two incompatible numbers without explanation.

\section{Multiplicity}

Test data was scored more than once during development, which weakens the
interpretation of the locked-test interval as a single pre-registered test. The
number of scorings is disclosed rather than minimised. The field evaluation was
performed once per checkpoint.

\section{What is not claimed}

No clinical or agronomic claim. No field-readiness claim. No deployment claim. The
demonstration interface replays recorded model output and is not a service. The
system is not a diagnostic instrument and must not be used as one.

No severity grading, no multi-disease frames, no species outside the label set.
